\metatitle{Positroids}
\metadescription{Background on positroids, Grassmann necklaces, planar networks, and positroid polytopes.}


\section[positroid]{Positroids}

A \defin{positroid}\label{positroidDefinition} is a
\hyperref[matroidDefinition]{matroid} that is represented by a point in the
totally nonnegative Grassmannian.  Positroids were introduced by
\name{Alexander Postnikov} in his study of
\hyperref[totalPositivity]{total positivity},
\hyperref[grassmannianVarieties]{Grassmannians}, and planar networks
\cite{Postnikov2006x}.  For the matrix total-positivity background used in path
models, see \hyperref[totallyNonnegative]{totally nonnegative matrices} and
\hyperref[tnnPathMatrix]{path matrices}.

The cyclic order on the ground set is part of the structure, so one usually
works with matroids on $[n]$ with its standard cyclic order
\[
  1,2,\dotsc,n,1,\dotsc .
\]
An arbitrary relabeling of a positroid need not preserve the positroid
property.

Let $A$ be a real $k\times n$ matrix of rank $k$, with columns labeled by
$[n]$.  For $I\in \binom{[n]}{k}$, let $\Delta_I(A)$ be the maximal minor using
the columns indexed by $I$.  The row span of $A$ is a point in the Grassmannian
$\mathrm{Gr}(k,n)$, and it is \defin{totally nonnegative} if $A$ can be chosen
so that
\[
  \Delta_I(A) \geq 0
  \qquad \text{for all } I\in \binom{[n]}{k}.
\]
The nonzero \hyperref[pluckerCoordinate]{Pl{\"u}cker coordinates} then define
a matroid:
\[
  \mathcal{B}(A)
  =
  \{ I\in \binom{[n]}{k} : \Delta_I(A) > 0 \}.
\]
A rank $k$ matroid on $[n]$ is a \defin{positroid} if it is of the form
$([n],\mathcal{B}(A))$ for such a totally nonnegative matrix $A$.

Equivalently, positroids index the cells in the totally nonnegative
Grassmannian:
\[
  \Pi_M =
  \left\{
    X\in \mathrm{Gr}(k,n)_{\geq 0} :
    \Delta_I(X) > 0 \text{ if and only if } I\in \mathcal{B}(M)
  \right\}.
\]
Postnikov proved that these cells are nonempty precisely for positroids
\cite{Postnikov2006x}.


\subsection[positroidGrassmannNecklace]{Grassmann necklaces}

There are several useful combinatorial encodings of positroids, including
\defin{decorated permutations}, Grassmann necklaces, Le-diagrams, and reduced
\defin{plabic graphs}\label{plabicGraphDefinition}.  The Grassmann-necklace
description is close to the language of
\hyperref[schubertMatroidDefinition]{Schubert matroids}.

For $a\in [n]$, let $\leq_a$ be the cyclic order
\[
  a \lt_a a+1 \lt_a \dotsb \lt_a n \lt_a 1 \lt_a \dotsb \lt_a a-1.
\]
If $I=\{i_1\lt_a \dotsb \lt_a i_k\}$ and
$J=\{j_1\lt_a \dotsb \lt_a j_k\}$, write $I\leq_a J$ if
$i_s\leq_a j_s$ for every $s$.

A \defin{Grassmann necklace}\label{grassmannNecklaceDefinition} of type
$(k,n)$ is a sequence
\[
  \mathcal{I}=(I_1,I_2,\dotsc,I_n),
  \qquad I_a\in \binom{[n]}{k},
\]
such that, with indices taken modulo $n$,
\[
  I_{a+1}=I_a \quad \text{if } a\notin I_a,
\]
and
\[
  I_{a+1}=(I_a\setminus \{a\})\cup \{b\}
  \quad \text{for some } b\in [n] \text{ if } a\in I_a.
\]

Given a Grassmann necklace $\mathcal{I}$, define
\[
  \mathcal{B}(\mathcal{I})
  =
  \left\{
    J\in \binom{[n]}{k} :
    I_a \leq_a J \text{ for all } a\in [n]
  \right\}.
\]
The condition $I_a\leq_a J$ says that $J$ is a basis in a
\hyperref[schubertMatroidDefinition]{Schubert matroid} with respect to the
cyclic order $\leq_a$.  Hence
$\mathcal{B}(\mathcal{I})$ is an intersection of cyclically shifted Schubert
conditions.

\begin{theorem}[Postnikov; Oh, \cite{Postnikov2006x,Oh2011}]
The collection $\mathcal{B}(\mathcal{I})$ is the set of bases of a positroid.
Conversely, every positroid arises uniquely in this way from its Grassmann
necklace.
\end{theorem}

For a positroid $M$, the set $I_a$ in the corresponding Grassmann necklace is
the lexicographically minimal basis of $M$ with respect to the cyclic order
$\leq_a$.


\subsection[positroidNetworks]{Planar networks and plabic graphs}

Planar networks give a concrete source of totally nonnegative matrices.
Consider an acyclic directed planar network with boundary vertices labeled by
$[n]$, and put nonnegative weights on the edges.  A
\defin{boundary measurement matrix}\label{boundaryMeasurementMatrix} records
weighted sums of paths from boundary sources to boundary sinks.  By the
\hyperref[lgvLemma]{Lindstr\"om--Gessel--Viennot lemma}, its maximal minors are
sums of weights of \hyperref[nonintersectingPaths]{nonintersecting} path
families, and hence are nonnegative.
The sets of boundary vertices that occur as nonzero maximal minors form a
positroid.  This is the same path-matrix mechanism that appears in
\hyperref[tnnPathMatrix]{total-positivity arguments}.

Postnikov showed that every positroid can be obtained from such planar
network data, equivalently from a reduced plabic graph
\cite{Postnikov2006x}.  This is one reason positroids occur naturally in
families built from lattice paths and planar placements.


\section[positroidExamples]{Examples of positroids}

\begin{example}[Uniform matroids]
Every \hyperref[matroidUniform]{uniform matroid} $U_{k,n}$ is a positroid.
Choose real numbers
$t_1\lt t_2\lt \dotsb \lt t_n$ and let
\[
  A =
  \begin{bmatrix}
    1 & 1 & \dotsb & 1 \\
    t_1 & t_2 & \dotsb & t_n \\
    \vdots & \vdots & & \vdots \\
    t_1^{k-1} & t_2^{k-1} & \dotsb & t_n^{k-1}
  \end{bmatrix}.
\]
For $I=\{i_1\lt \dotsb \lt i_k\}$, the minor $\Delta_I(A)$ is the Vandermonde
determinant
\[
  \prod_{1\leq p\lt q\leq k} (t_{i_q}-t_{i_p}),
\]
which is positive.  Thus every $k$-subset of $[n]$ is a basis.
\end{example}

\begin{example}[Schubert matroids]
Let $I=\{i_1\lt \dotsb \lt i_k\}\subseteq [n]$.  The corresponding
\hyperref[schubertMatroidDefinition]{Schubert matroid} has bases
\[
  \{j_1\lt \dotsb \lt j_k\}
  \qquad \text{with} \qquad
  i_s\leq j_s \text{ for all } s.
\]
\hyperref[schubertMatroidDefinition]{Schubert matroids}, and their cyclic
shifts, are positroids.  More generally,
positroids are exactly the compatible intersections of these cyclic Schubert
matroids described by Grassmann necklaces \cite{Oh2011}.
\end{example}

\begin{example}[Lattice path matroids]
Every \hyperref[latticePathMatroidDefinition]{lattice path matroid} is a
positroid \cite{Oh2011}.  In particular, the Catalan matroids obtained from
Dyck paths are positroids.  For the definition and examples, see
\hyperref[latticePathMatroids]{the page on lattice path matroids}.
\end{example}

\begin{example}[Rook matroids]
Alexandersson and Jal proved that every rook matroid is a positroid
\cite{AlexanderssonJal2024x}.  For the construction from
\hyperref[nonNestingRookPlacement]{nonnesting rook placements} on skew Ferrers
boards, see \hyperref[rookMatroids]{the rook matroids section}.
\end{example}


\section[positroidProperties]{Basic properties}

Positroids inherit many of the usual features of representable matroids, but
with extra compatibility with the cyclic order on $[n]$.

\subsection[positroidOperations]{Operations preserving positroids}

Positroids are closed under cyclic shifts of the ground-set order,
\hyperref[matroidDual]{matroid duality}, \hyperref[matroidRestriction]{restriction},
\hyperref[matroidDeletion]{deletion}, and
\hyperref[matroidContraction]{contraction}.  In particular, every
\hyperref[matroidMinors]{minor} of a
positroid is again a positroid \cite{ArdilaRinconWilliams2014}.

In the Grassmannian picture, duality corresponds to passing from a subspace
to its orthogonal complement, while deletion and contraction correspond to
removing or forcing a boundary element.

Direct sums require some compatibility with the cyclic order.  If the summands
occupy cyclic intervals, or more generally the blocks of a non-crossing
partition of the ground set, then the direct sum of positroids is again a
positroid \cite{ArdilaRinconWilliams2014}.


\subsection[positroidPolytopes]{Positroid polytopes}

Let $M$ be a rank $k$ positroid on $[n]$.  Its
\hyperref[matroidBasePolytope]{matroid base polytope} is
\[
  P_M = \operatorname{conv}\{\mathbf{e}_B : B\in \mathcal{B}(M)\}
  \subseteq \setR^n.
\]
For a \defin{cyclic interval}\label{cyclicIntervalDefinition} $[i,j]$, write
\[
  x_{[i,j]} \coloneqq \sum_{a\in [i,j]} x_a.
\]
Then $P_M$ can be described by the hypersimplex equation
\[
  x_1+x_2+\dotsb+x_n = k,
\]
the inequalities $x_i\geq 0$, and inequalities of the form
\[
  x_{[i,j]} \leq r_M([i,j])
\]
for cyclic intervals $[i,j]\subseteq [n]$.  Conversely, a matroid on $[n]$ is
a positroid if and only if its base polytope admits such a cyclic-interval
description \cite{ArdilaRinconWilliams2014}.

Thus \defin{positroid polytopes}\label{positroidPolytopeDefinition} form a
distinguished subclass of matroid base polytopes with far fewer defining rank
inequalities than a general matroid base polytope.  Moreover, every face of a
positroid polytope is again a positroid polytope, and the face poset embeds
naturally into a poset of weighted noncrossing partitions
\cite{ArdilaRinconWilliams2014}.

The closure of a positroid cell in the ordinary Grassmannian is called a
\defin{positroid variety}\label{positroidVarietyDefinition}.  These varieties
include \hyperref[schubertVarietyGeometryDefinition]{Schubert varieties} and
are one of the geometric reasons that Schubert matroids and their cyclic shifts
appear in the combinatorics of positroids.
\name{Thomas Lam} shows that graded pieces of the homogeneous coordinate ring
of a positroid variety are intersections of cyclically rotated rectangular
Demazure modules \cite{Lam2019}.  These cyclic Demazure modules have canonical
bases and cyclic Demazure crystals.
